\section{Expériences de récupération}\label{sec:recup}

Toutes les variantes renvoient six articles et sont comparées avec les mêmes métriques. \textbf{A1} est v1.1 avec l'indice
de code du jeu de données (une information dont un usager ne dispose pas) ; \textbf{A2} est v1.1 sans cet indice, la référence
équitable de l'agent.

\subsection{Planification contre question brute : un résultat négatif}
La première idée était de faire écrire par le LLM des requêtes de recherche (nœud Plan) au lieu de chercher avec la question
de l'usager. Sur le développement, plusieurs variantes de planification dépassent A2 : la variante D (requêtes du Plan,
code en filtre doux) obtient \NPvrDevDHitsix~\% contre \NPvrDevATwoHitsix~\% pour A2. Sur le test, l'avantage disparaît :
D obtient \NPvrTestDHitsix~\% et A2 \NPvrTestATwoHitsix~\% (tableau~\ref{tab:plan}, figure~\ref{fig:devtest}). Toutes les
variantes de planification baissent entre le développement et le test, alors que A2 monte. Seule B2 (requêtes du Plan, sans
filtre) reste au-dessus de A2 en Hit@6 sur le test (\NPvrTestBTwoHitsix~\% contre \NPvrTestATwoHitsix~\%), mais en dessous en
Hit@1 (\NPvrTestBTwoHitone~\% contre \NPvrTestATwoHitone~\%) et en MRR (\NPvrTestBTwoMrr{} contre \NPvrTestATwoMrr). Nous
concluons que la planification systématique ne bat pas nettement la question brute, et que l'avantage observé sur le
développement venait en partie de l'ajustement à ces \NQDev{} questions. Le Plan a été retiré du graphe.

\begin{table}[H]
\centering
\small
\begin{tabular}{lrrrr}
\toprule
Variante & Hit@6 dev & MRR dev & Hit@6 test & MRR test \\
\midrule
A1 & 88{,}6 & 0{,}659 & 80{,}8 & 0{,}687 \\
A2 & 71{,}4 & 0{,}544 & 75{,}0 & 0{,}646 \\
B1 & 80{,}0 & 0{,}622 & 73{,}1 & 0{,}610 \\
B2 & 85{,}7 & 0{,}592 & 76{,}9 & 0{,}630 \\
C & 80{,}0 & 0{,}611 & 73{,}1 & 0{,}635 \\
D & 77{,}1 & 0{,}630 & 73{,}1 & 0{,}609 \\
E & 80{,}0 & 0{,}607 & 73{,}1 & 0{,}638 \\
G & 80{,}0 & 0{,}650 & -- & -- \\
\bottomrule
\end{tabular}
\caption{Planification contre question brute (dev : 35 questions, test : 52). A1 utilise l'indice de code du jeu de données ; A2 est la référence équitable. G n'a été évaluée que sur le dev.}
\label{tab:plan}
\end{table}

\figfile{f3_dev_vs_test}{Hit@6 de A2 et de la variante D sur le développement puis sur le test : l'avantage de D sur le
développement ne se retrouve pas sur le test.}{fig:devtest}

\subsection{Indice de code doux (G)}
Plutôt que de remplacer la question, la variante G garde la recherche de A2 et lui ajoute, par fusion, la même recherche
filtrée sur le code proposé par l'Analyse lorsqu'il est unique. Sur le développement, G obtient \NLoopGHitsix~\% de Hit@6
contre \NLoopAtwoHitsix~\% pour A2. Les descriptions de codes utilisées par l'Analyse ont été écrites après avoir vu des
échecs du développement : la précision de l'indice y est donc optimiste.

\subsection{Boucle Évaluer → Reformuler}
La boucle part d'une récupération (A2 ou G), demande au LLM si les preuves suffisent et, sinon, reformule et cherche à
nouveau, au plus deux fois. L'attente écrite avant l'expérience était un gain d'au moins 3 points de Hit@6 sur A2, sans
concept dégradé. Partant de G, la boucle atteint \NLoopLoopGHitsix~\% ; partant de A2, \NLoopLoopAtwoHitsix~\%.
Partant de G, \NLoopGRetried{} questions ont été reprises : \NLoopGBetter{} se sont améliorées, \NLoopGSame{} sont restées
identiques et \NLoopGWorse{} se sont dégradées. Partant de A2 : \NLoopAtwoRetried{} reprises, \NLoopAtwoBetter{} améliorées,
\NLoopAtwoSame{} identiques, \NLoopAtwoWorse{} dégradées.

\subsection{Ablation et incertitude}
Le tableau~\ref{tab:ablation} et la figure~\ref{fig:ablation} placent les variantes entre A2 et l'oracle A1. L'écart
entre la boucle partant de G et A2 est de \NLoopLoopGDiff{} points, avec un intervalle de confiance à 95~\% de
[\NLoopLoopGDifflo{} ; \NLoopLoopGDiffhi] points en rééchantillonnant les concepts. L'intervalle est large : les gains
viennent de peu de concepts, et le développement a déjà servi à choisir ces variantes. Ce résultat reste donc
\emph{non confirmé} ; le jeu de test ayant déjà été vu, la confirmation viendra du nouveau jeu de cent questions.

\begin{table}[H]
\centering
\small
\begin{tabular}{lrrrr}
\toprule
Variante (dev) & Hit@6 & IC 95 \% & Écart vs A2 & IC 95 \% de l'écart \\
\midrule
A2 & 71{,}4 & [46{,}9 ; 91{,}7] & -- & -- \\
G & 80{,}0 & [57{,}7 ; 96{,}9] & $+$8{,}6 & [0{,}0 ; 23{,}5] \\
Boucle depuis A2 & 80{,}0 & [60{,}7 ; 95{,}5] & $+$8{,}6 & [0{,}0 ; 23{,}1] \\
Boucle depuis G & 85{,}7 & [68{,}8 ; 100{,}0] & $+$14{,}3 & [2{,}9 ; 30{,}8] \\
A1 (oracle) & 88{,}6 & [74{,}3 ; 100{,}0] & $+$17{,}1 & [3{,}1 ; 34{,}6] \\
\bottomrule
\end{tabular}
\caption{Ablation sur le dev (35 questions, 15 concepts). Intervalles par bootstrap apparié rééchantillonnant les concepts (5000 tirages, graine 20261022).}
\label{tab:ablation}
\end{table}

\figfile{f4_ablation}{Ablation sur le développement : de A2 à l'oracle A1. Barres d'erreur : intervalles à 95~\% par
bootstrap sur les concepts.}{fig:ablation}

\subsection{Par langue}
Le tableau~\ref{tab:bylang} et la figure~\ref{fig:bylang} détaillent Hit@6 par langue. Les effectifs par langue sont très
petits sur le développement ; nous ne tirons pas de conclusion sur un avantage d'une langue.

\begin{table}[H]
\centering
\small
\begin{tabular}{lrrr}
\toprule
Variante (dev) & FR (n=19) & AR (n=8) & EN (n=8) \\
\midrule
A2 & 73{,}7 & 62{,}5 & 75{,}0 \\
G & 84{,}2 & 62{,}5 & 87{,}5 \\
Boucle depuis A2 & 78{,}9 & 75{,}0 & 87{,}5 \\
Boucle depuis G & 89{,}5 & 75{,}0 & 87{,}5 \\
A1 (oracle) & 89{,}5 & 75{,}0 & 100{,}0 \\
\bottomrule
\end{tabular}
\caption{Hit@6 par langue sur le dev.}
\label{tab:bylang}
\end{table}

\figfile{f5_by_language}{Hit@6 par langue sur le développement, pour chaque variante.}{fig:bylang}

\subsection{Qualité du jugement « preuves suffisantes »}
La figure~\ref{fig:grade} croise le premier verdict du nœud Évaluer avec la présence réelle de l'article de référence
au tour 0. Partant de A2, le nœud a jugé suffisantes \NGradeAtwoEnoughMiss{} récupérations où l'article de référence
manquait, et faibles \NGradeAtwoWeakHit{} récupérations qui le contenaient ; partant de G, il a jugé
suffisantes \NGradeGEnoughMiss{} récupérations sans l'article de référence et faibles \NGradeGWeakHit{} récupérations
qui le contenaient. Les « suffisant » à tort sont des questions que la boucle ne peut pas rattraper.

\figfile{f6_grade_confusion}{Premier verdict du nœud Évaluer contre la présence de l'article de référence au tour 0,
pour les deux points de départ.}{fig:grade}

\subsection{Coût}
La boucle coûte en moyenne \NLoopGLlmper{} appels LLM (Évaluer et Reformuler) par question en partant de G, et
\NLoopAtwoLlmper{} en partant de A2 (figure~\ref{fig:reprises}).

\figfile{f7_retries_calls}{Nombre de reprises et nombre d'appels LLM de la boucle par question.}{fig:reprises}
